\SourceLeaf{352}%
kel eines geradlinigen Dreiecks \(4''95113,\ 4''95104,\ 4''95131\). Übrigens hat der Verf. auch die in den obigen Ausdrücken fehlenden Glieder der vierten Ordnung entwickelt, die für die Kugelfläche eine sehr einfache Form erhalten; bei messbaren Dreiecken auf der Oberfläche der Erde sind sie aber ganz unmerklich, und in dem angeführten Beispiel würden sie die erste Reduction nur um zwei Einheiten der fünften Decimale vermindert und die dritte eben so viel vergrössert haben.

\pdfbookmark[0]{Anzeigen. Untersuchungen über Gegenstände der höhern Geodaesie. 1843.}{anzeigen-geodaesie-erste}
\begin{center}
\rule{.85\textwidth}{.3pt}\par
Göttingische gelehrte Anzeigen.\quad 1843 November 6.\par
\rule{.85\textwidth}{.3pt}
\end{center}
\hspace*{\parindent}Der königlichen Societät ist am 23. October von dem Hofrath \textsc{Gauss} eine Vorlesung überreicht, mit der Überschrift:
\begin{center}\textit{Untersuchungen über Gegenstände der höhern Geodaesie},\end{center}
\noindent von welcher hier ein kurzer Bericht gegeben werden soll.

Bei dem trigonometrischen Theile der von dem Verf. in den Jahren 1821---1827 ausgeführten Gradmessung, und bei den spätern damit zusammenhängenden und über das ganze Königreich Hannover sich erstreckenden trigonometrischen Vermessungen sind, sowohl in Beziehung auf die Art, wie die Messungen angestellt wurden, als noch mehr in Beziehung auf ihre nachherige mathematische Behandlung und ihre Verarbeitung zu Resultaten, Wege eingeschlagen, die von den sonst betretenen abwichen. Manches von diesen dem Hofr. \textsc{Gauss} eigenthümlichen Methoden ist zwar bereits zur Öffentlichkeit gebracht, theils von ihm selbst in verschiedenen vorlängst erschienenen Aufsätzen, theils durch andere, welche nach mündlichen oder brieflichen Mittheilungen bei ihren eigenen trigonometrischen Messungen Anwendungen davon gemacht hatten. Allein der erheblichere Theil jener Methoden, diejenigen, welche sich am meisten von den sonst gebräuchlichen unterscheiden, und deren Verständniss eine tiefere mathematische Begründung erfordert, ist bisher noch nicht dargestellt. Des Verf. frühern Vorsatz, nach völliger Beendigung der Messungen diese selbst nebst allen von ihm angewandten Verfahrungsarten in einem besondern Werke darzulegen, haben Umstände, deren Auseinandersetzung nicht hieher gehört, zur Zeit noch procrastinirt, und er hat deshalb das Auskunftsmittel gewählt, das im theoreti%
